\section{Related Work and Positioning}

\paragraph{Sparse MoE routing and load balance.}
Sparse MoE models route each token to a small number of experts, increasing
total parameter capacity without activating all parameters for every token
\citep{shazeer2017outrageously,fedus2022switch}. This makes the router a
central object: it controls not only inference-time computation, but also which
experts receive training signal. Practical MoE systems often introduce
load-balancing losses, capacity factors, token dropping, or assignment rules to
avoid severe expert underuse. Those mechanisms are important for engineering
stable training, but they also make interpretation harder. If a count-balancing
mechanism is active, uniform expert usage may reflect the constraint rather
than a naturally balanced feature geometry. The present paper isolates the
unconstrained top-1 case so that this interpretive issue is visible.

\paragraph{Geometry before semantic expertise.}
A route pattern is sometimes tempting to read as evidence that different
experts have learned different semantic functions. This paper uses a narrower
standard. A top-1 linear router first sees inner-product geometry. Therefore a
balanced or imbalanced load vector can be caused by the angular distribution of
hidden states and by the geometry of router rows before any claim about expert
utility is justified. The controlled spherical model is a way to state this
minimal geometry exactly. It does not deny that semantic specialization can
exist; it says that semantic specialization requires evidence beyond count
balance, such as route-function binding, expert-specific utility, or a causal
intervention on expert behavior.

\paragraph{Spectral router and expert geometry.}
Recent router designs and analyses motivate using expert matrices themselves
to define routing directions. A manifold power-iteration style update, for
example, can align a router row with an expert's principal input-side direction
\citep{routermpi2026}. This is a meaningful way to construct candidate centers.
However, center construction is not identical to cell balance. If the recovered
expert directions cluster, then the induced top-1 cells can have unequal mass
even when each row is well aligned with its expert. This paper therefore treats
spectral alignment as a proxy-construction question and load balance as a
separate cell-mass question.

\paragraph{Controlled baselines for ongoing training studies.}
The broader research line motivating this manuscript studies early training
dynamics, common hidden-state components, and preservation of routing
partitions. Those questions require additional objects: margins, drift budgets,
loss guards, and route-function utility. The present paper deliberately stops
before those dynamic variables. Its role is to define the static baseline that
the dynamic analysis must respect. If a future training-time method claims to
produce meaningful uniform feature routing, it should first pass the static
checks developed here and then add evidence that the partition survives
optimization and supports useful expert computation.

